\documentclass[10pt]{article}
% Shared style for BELAND-PLUS academic revision R3
% Typography: Latin Modern (the canonical scalable Computer Modern family).
\usepackage[T1]{fontenc}
\usepackage[utf8]{inputenc}
\usepackage{lmodern}
\usepackage{microtype}
\usepackage{amsmath,amssymb,mathtools,bm}
\usepackage{booktabs,threeparttable,longtable,tabularx,array,multirow,makecell}
\usepackage{graphicx,adjustbox}
\usepackage{caption,subcaption}
\usepackage{tikz}
\usetikzlibrary{arrows.meta,positioning,fit,calc,backgrounds,matrix,decorations.pathreplacing,shapes.geometric}
\usepackage{pgfplots}
\pgfplotsset{compat=1.18}
\usepackage{xcolor}
\usepackage{geometry}
\geometry{margin=1in,headheight=15pt,footskip=30pt}
\usepackage{setspace}
\setstretch{1.10}
\usepackage{enumitem}
\setlist[itemize]{leftmargin=1.55em,itemsep=0.18em,topsep=0.35em,parsep=0pt}
\setlist[enumerate]{leftmargin=1.8em,itemsep=0.18em,topsep=0.35em,parsep=0pt}
\usepackage{titlesec}
\titleformat{\section}{\large\bfseries}{\thesection}{0.75em}{}
\titleformat{\subsection}{\normalsize\bfseries}{\thesubsection}{0.65em}{}
\titleformat{\subsubsection}{\normalsize\itshape}{\thesubsubsection}{0.55em}{}
\titlespacing*{\section}{0pt}{2.0ex plus .4ex minus .2ex}{0.9ex}
\titlespacing*{\subsection}{0pt}{1.55ex plus .3ex minus .2ex}{0.65ex}
\titlespacing*{\subsubsection}{0pt}{1.1ex plus .2ex minus .2ex}{0.45ex}
\usepackage{fancyhdr}
\pagestyle{fancy}
\fancyhf{}
\fancyhead[L]{\footnotesize\itshape Hurricane Ida and public emergency-service observability}
\fancyhead[R]{\footnotesize BELAND-PLUS}
\fancyfoot[C]{\thepage}
\renewcommand{\headrulewidth}{0.35pt}
\renewcommand{\footrulewidth}{0pt}
\usepackage[round,authoryear]{natbib}
\usepackage{hyperref}
\usepackage[nameinlink,noabbrev]{cleveref}
\usepackage{xurl}
\urlstyle{same}
\usepackage{csquotes}
\usepackage{listings}
\usepackage{tcolorbox}
\tcbuselibrary{breakable,skins}
\usepackage{float}
\usepackage[section]{placeins}
\usepackage{siunitx}
\sisetup{detect-all,group-separator={,},group-minimum-digits=4,table-number-alignment=center}
\usepackage{pdflscape}
\usepackage{etoolbox}

\definecolor{BelandBlue}{HTML}{183A5A}
\definecolor{BelandBlueLight}{HTML}{EAF0F5}
\definecolor{BelandOrange}{HTML}{B5652A}
\definecolor{BelandRed}{HTML}{9E3D36}
\definecolor{BelandGray}{HTML}{4B5563}
\definecolor{BelandLightGray}{HTML}{F4F5F7}
\definecolor{BelandGreen}{HTML}{2E6F64}
\definecolor{BelandGold}{HTML}{A47A24}

\hypersetup{
  colorlinks=true,
  linkcolor=BelandBlue,
  citecolor=BelandBlue,
  urlcolor=BelandBlue,
  pdfborder={0 0 0}
}

\newcommand{\Jzerozero}{J_{00}}
\newcommand{\Jzeroone}{J_{01}}
\newcommand{\Jonezero}{J_{10}}
\newcommand{\Joneone}{J_{11}}
\newcommand{\E}{\mathrm{E}}
\newcommand{\R}{\mathrm{R}}
\newcommand{\Prob}{\Pr}
\newcommand{\Ida}{\mathrm{Ida}}
\newcommand{\Iset}{\mathcal{I}}
\newcommand{\Mset}{\mathfrak{M}}
\newcommand{\Xset}{\mathcal{X}}
\newcommand{\supp}{\operatorname{supp}}
\newcommand{\diam}{\operatorname{diam}}
\newcommand{\argmin}{\operatorname*{arg\,min}}
\newcommand{\argmax}{\operatorname*{arg\,max}}
\newcommand{\ind}{\mathbf{1}}
\newcommand{\Rset}{\mathbb{R}}
\newcommand{\Nset}{\mathbb{N}}
\newcommand{\abs}[1]{\left|#1\right|}
\newcommand{\norm}[1]{\left\lVert#1\right\rVert}
\newcommand{\given}{\,|\,}

\newtcolorbox{plainbox}[1][]{
  enhanced,breakable,colback=BelandBlueLight,colframe=BelandBlue,
  boxrule=0.55pt,arc=1.5mm,left=2.7mm,right=2.7mm,top=1.8mm,bottom=1.8mm,#1}
\newtcolorbox{resultbox}[1][]{
  enhanced,breakable,colback=BelandLightGray,colframe=BelandGray,
  boxrule=0.5pt,arc=1.5mm,left=2.7mm,right=2.7mm,top=1.8mm,bottom=1.8mm,#1}
\newtcolorbox{warningbox}[1][]{
  enhanced,breakable,colback=white,colframe=BelandOrange,
  boxrule=0.65pt,arc=1.5mm,left=2.7mm,right=2.7mm,top=1.8mm,bottom=1.8mm,#1}
\newtcolorbox{definitionbox}[1][]{
  enhanced,breakable,colback=white,colframe=BelandBlue!70,
  boxrule=0.45pt,arc=1.2mm,left=2.5mm,right=2.5mm,top=1.5mm,bottom=1.5mm,#1}

\captionsetup{font=small,labelfont=bf,labelsep=period,skip=5pt,justification=raggedright,singlelinecheck=false}
\renewcommand{\arraystretch}{1.18}
\setlength{\tabcolsep}{5pt}
\setlength{\parindent}{1.35em}
\setlength{\parskip}{0pt}
\setlength{\emergencystretch}{2em}
\raggedbottom

\lstdefinestyle{belandcode}{
  basicstyle=\ttfamily\footnotesize,
  backgroundcolor=\color{BelandLightGray},
  frame=single,
  rulecolor=\color{BelandGray!45},
  breaklines=true,
  breakatwhitespace=false,
  columns=fullflexible,
  keepspaces=true,
  showstringspaces=false,
  xleftmargin=0.5em,
  xrightmargin=0.5em,
  aboveskip=0.7em,
  belowskip=0.7em
}
\lstset{style=belandcode}

\newcommand{\notclaim}[1]{\textit{This does not identify #1.}}
\newcommand{\pp}{percentage points}
\newcommand{\keywords}[1]{\par\noindent\textbf{Keywords:} #1}
\newcommand{\jel}[1]{\par\noindent\textbf{JEL codes:} #1}

\AtBeginEnvironment{tabular}{\small}
\AtBeginEnvironment{tabularx}{\small}
\AtBeginEnvironment{longtable}{\small}

\geometry{margin=0.72in,headheight=13pt,footskip=22pt}
\setstretch{1.02}
\fancyhead[L]{\footnotesize\itshape Research status note}
\fancyhead[R]{\footnotesize 26 August 2026}

\title{\vspace{-2.2em}\textbf{Research Status Note}\\[0.25em]
\large Public CAD Is Not Operational Ground Truth --- R16.2}
\author{Yuwen Zhu\\Carleton University}
\date{26 August 2026}

\begin{document}
\maketitle
\thispagestyle{plain}
\vspace{-1.9em}

\begin{plainbox}[title=Four findings to retain]
The released public-field convention changes abruptly on 28 July 2021. Hurricane Ida produces a second, temporary change inside that later convention. Its largest half-day field-state change is 53.2 percentage points, larger than all 151 post-change ordinary-week comparisons. The public record does not reveal physical response or the institutional mechanism that generated the released fields.
\end{plainbox}

\noindent\textbf{Evidence behind the sequence.}
\begin{itemize}
\item In 419,840 non-officer-initiated records created through 27 July 2021, none has a valid arrival field without a dispatch field. The configuration first appears on 28 July, when officer-stream dispatch coverage also falls from 100 percent to 40 percent and then to 1 percent on 29 July.
\item During Ida, the public missing-dispatch share reaches 43.4 percent on 31 August and 49.3 percent on 1 September. The 29 August arrival-field dip and a smaller 10 September excursion show that missingness changes form across days.
\item The full-count post-change comparison is primary. A standardized comparison also places Ida first but covers only 86.1 percent of event and 77.0 percent of baseline records and fails the 0.90 reference eligibility rule.
\item The comparable non-officer dispatch-field series is 65.4 percent in 2024, 66.0 percent in 2025, and 66.8 percent in the 2026 snapshot. Pooling officer-initiated records produces a different denominator and a lower all-row rate.
\item Public documents define field labels and exclude two proposed technology explanations. They do not supply a versioned 2021 changelog or an internal-to-public lineage bridge.
\end{itemize}

\noindent\textbf{External relevance.} The paper now connects its public-record stress test to environmental-stress studies of call volume, police activity, priority groups, and response time. Its validation protocol asks researchers to declare call initiation, define the response clock, report endpoint coverage, and document priority and data-lineage rules before assigning an operational interpretation.

\noindent\textbf{Version boundary.} R16.2 broadens the paper's external relevance to environmental-stress studies without changing the validated R15.1 results. No numerical result was changed or promoted.

\noindent\textbf{Deliverables.} A standalone main paper, empirical supplement, this one-page note, reproduction manual, and separate legacy mathematical/formal-verification archive. The combined reviewer PDF contains only the main paper and empirical supplement.

\end{document}
